\documentclass[10.5pt]{article}
\usepackage{fontspec}\defaultfontfeatures{Ligatures=TeX}
\usepackage[letterpaper,margin=0.6in]{geometry}
\usepackage{amsmath,xcolor,longtable,array,tikz}
\usetikzlibrary{arrows.meta}
\setlength{\parindent}{0pt}\setlength{\parskip}{4pt}
\definecolor{pblue}{HTML}{1971C2}\definecolor{pred}{HTML}{E03131}\definecolor{pgreen}{HTML}{2F9E44}\definecolor{ppurple}{HTML}{7048E8}\definecolor{pgray}{HTML}{868E96}
\newfontfamily\mono{TeX Gyre Cursor}
\newcommand{\mtag}[2]{\colorbox{#1}{\textcolor{white}{\textsf{\fontsize{6.5}{7}\selectfont\bfseries #2}}}}
\newcommand{\clock}[1]{{\mono\fontsize{7}{8}\selectfont #1}\ }
% numbered step: \step{num}{tagcolor}{TAG}{text}
\newcommand{\step}[4]{\par\noindent\makebox[0.32in][l]{\textbf{\small #1}}\mtag{#2}{#3}\ {\small #4}\par\vspace{2pt}}
\newcommand{\pic}[1]{\par\hspace{0.32in}\begin{tikzpicture}[scale=0.42,every node/.style={font=\scriptsize}]#1\end{tikzpicture}\par\vspace{3pt}}
\newcommand{\CHECK}[3]{\par\noindent\makebox[0.32in][l]{\textbf{\small #1}}\fcolorbox{black}{black!4}{\parbox{2.55in}{\textsf{\fontsize{6.5}{8}\selectfont\bfseries CHECK}\ {\fontsize{7.5}{9}\selectfont #2}\par{\fontsize{7.5}{9}\selectfont\textbf{Answer:} #3}\par{\fontsize{7}{8.5}\selectfont\color{pgray} 4 of 5 right: go on. Fewer: one more example, ask again.}}}\par\vspace{3pt}}
\newcommand{\WATCH}[1]{\par\hspace{0.32in}{\textsf{\fontsize{6.5}{8}\selectfont\bfseries\color{pred} WATCH FOR}\ {\small\itshape\color{pred} #1}}\par}
\newcommand{\bl}[1]{\underline{\textcolor{pgreen}{\textbf{#1}}}}
\newcommand{\pen}[2]{\textcolor{#1}{\rule{8pt}{8pt}}\ #2}
\newcommand{\picL}[1]{\par\hspace{0.32in}\begin{tikzpicture}[scale=0.6,every node/.style={font=\footnotesize}]#1\end{tikzpicture}\par\vspace{3pt}}
% shell picture: x, y, shells (list of electron counts is drawn by \shells), label
\newcommand{\ring}[4]{\draw[pblue] (#1,#2) circle (#3);\foreach \k in {1,...,#4}{\pgfmathsetmacro{\ang}{90+360*\k/#4}\fill ({#1+#3*cos(\ang)},{#2+#3*sin(\ang)}) circle (0.07);}}
\newcommand{\nuc}[3]{\fill[pred!20] (#1,#2) circle (0.3);\node[pred] at (#1,#2){\tiny #3};}
% Lewis dots: \sd{x}{y}{n}{offset x}{offset y}{pair spread x}{pair spread y}
\newcommand{\sd}[7]{\ifnum#3=1 \fill ({#1+#4},{#2+#5}) circle (0.07);\fi\ifnum#3=2 \fill ({#1+#4-#6},{#2+#5-#7}) circle (0.07);\fill ({#1+#4+#6},{#2+#5+#7}) circle (0.07);\fi}
% \lw{x}{y}{symbol}{top}{right}{bottom}{left}: a box with the symbol and the dots
\newcommand{\lw}[7]{\draw[pgray] (#1-0.65,#2-0.65) rectangle (#1+0.65,#2+0.65);\node at (#1,#2){\scriptsize #3};\sd{#1}{#2}{#4}{0}{0.45}{0.14}{0}\sd{#1}{#2}{#5}{0.45}{0}{0}{0.14}\sd{#1}{#2}{#6}{0}{-0.45}{0.14}{0}\sd{#1}{#2}{#7}{-0.45}{0}{0}{0.14}}
\newcommand{\nb}[3]{\draw[pgray] (#1-0.65,#2-0.65) rectangle (#1+0.65,#2+0.65);\node at (#1,#2){\scriptsize #3};}
% bracket pair with symbol and charge: \bk{x}{y}{symbol}{charge}
\newcommand{\bk}[4]{\draw (#1-0.7,#2-0.65)--(#1-0.85,#2-0.65)--(#1-0.85,#2+0.65)--(#1-0.7,#2+0.65);\draw (#1+0.7,#2-0.65)--(#1+0.85,#2-0.65)--(#1+0.85,#2+0.65)--(#1+0.7,#2+0.65);\node at (#1,#2){\scriptsize #3};\node[pred,font=\tiny] at (#1+1.15,#2+0.8){#4};}
\pagestyle{plain}
\begin{document}
\begin{center}{\Large\bfseries Unit 4, Day 8 board script: Naming Molecules and Bond Types}\\[2pt]{\normalsize College prep \textbullet{} 50 minutes \textbullet{} \texttt{chemistry/whiteboard/index.html\#lesson=naming}}\end{center}
\noindent\textbf{Pens.} \pen{pblue}{shells, distance} \quad \pen{pred}{protons, pull} \quad \pen{pgreen}{new words (labels)} \quad \pen{ppurple}{shared electrons} \quad \pen{black}{everything else}\par
\noindent\textbf{Steps.} Numbered by frame. Each DRAW step has a small picture of the finished board. Tick the green Words box as you write each label.\par
\noindent\textbf{Not today.} No prefixes above penta-, no polyatomic ions, no polar bonds, no Roman numerals, no common names except as a joke. Naming is by the rule only.\par\smallskip
\noindent\fcolorbox{black}{black!3}{\parbox{\dimexpr\linewidth-2\fboxsep}{\small
\textbf{Time checkpoints.} Frame 2 by \textbf{13:00}, frame 3 by \textbf{22:00}, frame 4 by \textbf{32:00}, exit check by \textbf{39:00}. Checks are never cut.\par
\textbf{If you are behind,} cut in this order: (1) the second table in frame 2 (name to formula) becomes two whiteboard questions; (2) Al and SrCl$_2$ in the frame 3 table; (3) exit question a; (4) the summary becomes homework.\par
\textbf{If you are ahead,} on whiteboards: \emph{Name $\text{PBr}_3$.} (Phosphorus tribromide.) \emph{Formula for carbon tetrafluoride.} ($\text{CF}_4$.) \emph{Bond in $\text{CaF}_2$, in $\text{NF}_3$, in Ag?} (Ionic, covalent, metallic.) \emph{Why do ionic names have no prefixes?} (The charges set the ratio.)}}\par\medskip
\begin{longtable}{@{}>{\setlength{\parskip}{5pt}}p{3.55in}@{\hspace{0.25in}}p{3.05in}@{}}
\multicolumn{2}{@{}l}{\large\bfseries Frame 0: question and warm-up \hfill{\mono\small 0:00--7:00}}\tabularnewline*[2pt]
\textbf{Essential question.} How do we name a molecular compound, and how do the elements in a substance tell us which kind of bond holds it together?\par
\textbf{Warm-up.} Pairs shared by two nitrogen atoms? Lone pairs on oxygen in water? Name Rb$_2$O. &
\step{0.1}{black}{SAY}{\clock{0:00}Warm-up is up as they walk in. At 3:00 remove the covers: 3, a triple bond; 2; rubidium oxide.}
\step{0.2}{black}{SAY}{\clock{5:00}Read the question. ``Yesterday we drew molecules. Today we name them.''} \tabularnewline[10pt]
\multicolumn{2}{@{}l}{\large\bfseries Frame 1: prefixes \hfill{\mono\small 7:00--14:00}}\tabularnewline*[2pt]
Name the first element by its element name and the second with \bl{-ide}. A \bl{prefix} says how many atoms. Leave \bl{mono-} off the \bl{first} element. A compound made only of nonmetals is a \bl{molecular compound}.\par\smallskip
\textbf{1.} A student names SO$_3$ ``monosulfur trioxide.'' What is wrong? &
\step{1.1}{pgreen}{DRAW}{Fill the prefix table: mono 1, di 2, tri 3, tetra 4, penta 5. Say them aloud with the class.}
\pic{\foreach \i/\p/\n in {0/mono/1,1/di/2,2/tri/3,3/tetra/4,4/penta/5}{\node[font=\tiny] at (\i*1.6,0){\p-};\node[pgreen,font=\tiny] at (\i*1.6,-0.5){\n};}}
\step{1.2}{black}{DRAW}{NF$_3$: nitrogen trifluoride. Then N$_2$O: dinitrogen monoxide. Circle the mono- on oxide: ``Mono is dropped only on the first element.''}
\pic{\node[right] at (0,0){NF$_3$: nitrogen \textbf{tri}fluoride};\node[right] at (0,-0.8){N$_2$O: \textbf{di}nitrogen \textbf{mon}oxide};}
\step{1.3}{pgreen}{WORDS}{Tick \textbf{molecular compound}, \textbf{prefix}, \textbf{-ide}, \textbf{mono-}, \textbf{first}.}
\CHECK{1.4}{``Monosulfur trioxide'' for SO$_3$?}{Mono- is never used on the first element: sulfur trioxide.}
\WATCH{Dropping mono- from the second element too (``dinitrogen oxide'').} \tabularnewline[10pt]
\multicolumn{2}{@{}l}{\large\bfseries Frame 2: name and formula \hfill{\mono\small 14:00--22:00}}\tabularnewline*[2pt]
A prefix gives the \bl{subscript} of the element right after it. A missing prefix on the first element means \bl{1}.\par\smallskip
\textbf{2.} A student writes N$_3$O$_2$ for dinitrogen trioxide. What is wrong? &
\step{2.1}{black}{DRAW}{Fill the Name row: carbon disulfide, dichlorine monoxide, phosphorus pentachloride, silicon tetrafluoride. Say: ``Read the subscript, say the prefix.''}
\step{2.2}{black}{DRAW}{Second table: dinitrogen trioxide $\to$ N$_2$O$_3$; sulfur dichloride $\to$ SCl$_2$ (no prefix on sulfur means one). Tick \textbf{subscript} and \textbf{1}.}
\pic{\node[right] at (0,0){dinitrogen trioxide: N\textbf{$_2$}O\textbf{$_3$}};\node[right] at (0,-0.8){sulfur dichloride: SCl\textbf{$_2$}};}
\CHECK{2.3}{N$_3$O$_2$ for dinitrogen trioxide?}{The prefix goes with the element after it: di- on N (N$_2$), tri- on O (O$_3$).}
\WATCH{Putting the prefix number on the wrong element.} \tabularnewline[10pt]
\multicolumn{2}{@{}l}{\large\bfseries Frame 3: three kinds of bonds \hfill{\mono\small 22:00--31:00}}\tabularnewline*[2pt]
Look at what the substance is made of: metal + nonmetal is \bl{ionic}, nonmetal + nonmetal is \bl{covalent}, a metal alone is \bl{metallic}. The \bl{metallic bond} holds the atoms of a metal together.\par\smallskip
\textbf{3.} A student says H$_2$S is ionic because it contains sulfur. What is wrong? &
\step{3.1}{black}{DRAW}{Fill the top table: ionic, electrons move; covalent, electrons shared (purple); metallic, electrons move freely.}
\pic{\node[right] at (0,0.5){metal + nonmetal: ionic, \textbf{move}};\node[ppurple,right] at (0,-0.1){nonmetal + nonmetal: covalent, \textbf{shared}};\node[right] at (0,-0.7){metal only: metallic, \textbf{move freely}};}
\step{3.2}{black}{DRAW}{Bottom table, students call each one: LiBr ionic, NH$_3$ covalent, Zn metallic, SrCl$_2$ ionic, Al metallic. Tick \textbf{ionic}, \textbf{covalent}, \textbf{metallic}, \textbf{metallic bond}.}
\CHECK{3.3}{H$_2$S is ionic because it contains sulfur?}{Hydrogen and sulfur are both nonmetals: covalent.}
\WATCH{Calling a bond ionic because one of the elements is ``reactive.'' Ask: is there a metal?} \tabularnewline[10pt]
\multicolumn{2}{@{}l}{\large\bfseries Frame 4: sea of electrons \hfill{\mono\small 31:00--38:00}}\tabularnewline*[2pt]
In a metal, each atom gives its valence electrons to a shared \bl{sea of electrons}; the atoms are left as positive ions. The electrons \bl{move}, so a metal carries current. The ions slide past each other and the sea stays around them, so a metal \bl{bends} without breaking.\par\smallskip
\textbf{4.} A student says aluminum foil carries current because its atoms flow along the foil. What is wrong? &
\step{4.1}{pred}{DRAW}{Between the ions draw many small dots (use a different pen for the electrons if you like). Say: ``The positive ions stay; the electrons are free.''}
\pic{\foreach \i in {0,1,2,3}{\draw (\i*1.2,0) circle (0.28);\node[font=\tiny] at (\i*1.2,0){+};}\foreach \x/\y in {0.6/0.4,1.8/-0.4,3/0.4,0.5/-0.5,2.4/0.45,3.6/-0.4}{\fill[pred] (\x,\y) circle (1.2pt);}}
\step{4.2}{pgreen}{WORDS}{Tick \textbf{sea of electrons}, \textbf{move}, \textbf{bends}. Show bending with a hand: ``Layers slide; the sea stays.''}
\CHECK{4.3}{Aluminum foil carries current because its atoms flow?}{The atoms stay put. The free electrons in the sea move.}
\WATCH{Saying the electrons ``break'' when a metal bends.} \tabularnewline[10pt]
\multicolumn{2}{@{}l}{\large\bfseries Frame 5: exit check and summary \hfill{\mono\small 38:00--50:00}}\tabularnewline*[2pt]
\textbf{a.} Name NI$_3$; formula for oxygen difluoride. \textbf{b.} Bond type: Mg, HI, KF. \textbf{c.} Why can a metal bend without breaking?\par
\textbf{Summary.} Answer the essential question in two or three sentences. Use the words ``prefix'' and ``sea.'' &
\step{5.1}{black}{SAY}{\clock{38:00}Alone, silent, 6 minutes. Remove the covers at 44:00: nitrogen triiodide, OF$_2$; metallic, covalent, ionic; the ions slide and the sea stays around them.}
\step{5.2}{black}{SAY}{\clock{44:00}Summary: silent writing, read two aloud. Listen for ``prefix'' and ``sea.''}
\step{5.3}{black}{SAY}{\clock{48:00}``Tomorrow is practice for today and yesterday: Lewis structures, names and bond types.'' Stop there.} \tabularnewline
\end{longtable}
\end{document}
